%% GENERATED by python/paper/build.py from results/ -- do not edit by hand.
%% Rebuild:  .venv\Scripts\python.exe python\paper\build.py --only USUKFRCalibration_vectorX
%% Source:   results/paper/usCalibrationSummary.csv
\begin{table}[!htb]
\centering
\begin{threeparttable}
\caption{Calibration, US, UK, and France (vector $X_i$ calibration)}
\label{table:US:Calib_vectorX}
\renewcommand{\arraystretch}{1.25}
\begin{tabularx}{\textwidth}{Y|YYY|p{6cm}}
\hline
& \multicolumn{3}{c|}{\textbf{Country}} & \\ \cline{2-4}
\multicolumn{1}{c|}{\textbf{Parameter}} & \textbf{US} & \textbf{UK} & \textbf{France} & \textbf{Target} \\ \hline
$\theta$ & 0.74 & 0.56 & 1.00 & Replacement rate dispersion \\
$\omega$ & 1.45 & 1.16 & 1.42 & $\tau^{US} = 14.4\%$, $\tau^{UK} = 11.9\%$, $\tau^{FR} = 21.3\%$ \\
$\beta$ & 0.76 & 0.76 & 0.76 & US: 30y interest rate; imposed on UK/FR \\
$X$ & 0.88 & 1.55 & 1.33 & Avg.\ workweek \\
$\nu_{2020}$ & 1.34 & 1.17 & 0.97 & 30-year gross population growth rates \\
$\eta_{H}/\eta_L$ & 3.73 & 3.87 & 2.44 & Relative productivity of high (H) to low (L) income groups \\
\hline
\end{tabularx}
\begin{tablenotes}
\footnotesize
\item \textit{Note:} $\rho = 1.0$. $X$ is the population-weighted mean of $X_i$; its level is the hours unit, pinned for France and the UK by targeting average hours relative to the US rather than in levels. $\beta$ is calibrated for the US and imposed on the other two. Vector-$X_i$ calibration: $X_i$ is identified from relative hours, which are data here, and the level of $\bar h$ is then not identified --- only its ratio to the baseline is. $\beta$, $\omega$, the tax rate and the savings rate are common to the two variants; what differs is anything defined through $\eta$ or $X$.
\end{tablenotes}
\end{threeparttable}
\end{table}
